% Generated draft: requires mathematical and source review.
\documentclass[11pt]{article}
\usepackage[margin=25mm]{geometry}
\usepackage{fontspec}
\setmainfont{Latin Modern Roman}
\usepackage{amsmath,amssymb,mathtools,unicode-math}
\setmathfont{Latin Modern Math}
\usepackage{enumitem,needspace,xurl,hyperref}
\setlength{\parindent}{0pt}
\setlength{\parskip}{5pt}
\setlength{\emergencystretch}{3em}
\newcommand{\N}{\mathbb{N}}
\newcommand{\Z}{\mathbb{Z}}
\newcommand{\Q}{\mathbb{Q}}
\newcommand{\R}{\mathbb{R}}
\newcommand{\C}{\mathbb{C}}
\newcommand{\F}{\mathbb{F}}
\newcommand{\A}{\mathbb{A}}
\newcommand{\PP}{\mathbb{P}}
\newcommand{\E}{\mathbb{E}}
\newcommand{\eps}{\varepsilon}
\newcommand{\dd}{\,\mathrm{d}}
\begin{document}
\section*{UnsolvedMath ID: 8800076}
\textbf{Source problem code:} catalogue entry
\subsection{Definitions and mathematical statement}
The problem asks for cryptographic constructions whose security follows by a rigorous reduction from worst-case computational hardness in algebraic number theory. In an ideal-lattice setting, number-field ideals are viewed as Euclidean lattices through field embeddings. One example of a candidate hard problem is approximate ideal shortest vector: given an ideal lattice $I$, find a nonzero $v\in I$ with
\[
 \|v\|\le\gamma(d)\lambda_1(I),
 \qquad d=[K:\Q],
\]
for a specified approximation function $\gamma$. The source does not prescribe this exact problem, norm, field family, or approximation factor.

The required worst-case-to-security connection must turn an efficient successful adversary against the construction, averaged over its generated keys or challenges, into an algorithm for the stated hard algebraic problem on arbitrary admissible instances. This distinguishes a worst-case assumption from merely assuming that sampled instances are hard. The primitive, adversary model, parameters, failure probabilities, reduction losses, and meaning of ``directly'' are not specified in the short statement.

The supplied summary offers ideal/cyclic knapsack and ring-LWE as affirmative directions. Ring-LWE concerns noisy linear relations over number-field rings with chosen modulus, secret and error distributions. Its exact worst-case reductions can depend on field and parameter conditions, approximation factors, search-versus-decision form, and classical-versus-quantum algorithms. Those hypotheses are absent from the record, so the broad source phrase ``worst-case Ideal-SVP hardness implies security'' is not promoted to an unrestricted theorem for every number field or cryptosystem.
\subsection{Short English statement}
Build cryptographic primitives with security justified by worst-case hardness of specified algebraic-number-theory problems. The reduction and its parameter and computational-model conditions are part of the required construction.

The supplied summary labels the intended program solved, citing Micciancio's cyclic/ideal knapsack and ring-LWE work by Lyubashevsky--Peikert--Regev and others. It provides no exact theorem hypotheses, so the solved assessment is retained at program level rather than as a blanket reduction statement. Current primary-literature audit remains required.
\subsection{Sources}
\begin{itemize}
\item Local UnsolvedMath registry, permanent ID 8800076, source code catalogue entry, ``Cryptography from worst-case algebraic-number-theory hardness,'' proposed by Daniele Micciancio. Dataset curation and metadata: Ulam AI and UnsolvedMath Contributors, CC BY 4.0.
\item IPAM, \emph{Workshop on Number Theory and Cryptography --- Open Problems} (2006), Attachment 6240, PDF page 40, schedule locator \url{https://www.ipam.ucla.edu/programs/workshops/workshop-i-number-theory-and-cryptography-open-problems/?tab=schedule}; pdf-and-legacy-ppt visual-audit extraction, accessed 2026-07-29.
\item Supplied summary attributions: Micciancio (2002); Lyubashevsky--Peikert--Regev (2010); Peikert--Rosen; Stehlé--Steinfeld--Tanaka--Xagawa. Full titles, publication details, and exact reduction conditions are absent.
\item Source provenance reports inspection of 17 PDFs and two legacy PPT decks among 19 surviving attachments, excluding repetitions and solved definitions; this is supplied provenance, not a new inspection claimed here.
\item Current primary-literature status audit and redistribution terms NEEDS\_REVIEW for the publicly accessible author/publisher copy, separately from the dataset licence.
\end{itemize}
\end{document}
